\documentclass[12pt]{article} % Тип документа (стаття) та розмір шрифту (12pt)
\usepackage[T2A]{fontenc}      % Кодування шрифтів для відображення кирилиці
\usepackage[utf8]{inputenc}    % Кодування вхідного файлу (стандарт для укр.мови)
\usepackage[ukrainian]{babel}  % Підключення словників та правил переносу укр.мови
\usepackage{amsmath, amssymb}  % Пакети для математичних формул та символів
\usepackage{array}             % Пакет для покращеної роботи з таблицями

\title{Практична робота: Математичні формули} 
\author{Гуцол Альона}                          
\date{\today}                                

\begin{document} % Початок тексту документа

\maketitle % Команда, яка фактично малює заголовок, автора та дату на сторінці

\section*{Квадратне рівняння}

Квадратне рівняння має вигляд: $ax^2 + bx + c = 0$

\begin{equation}
ax^2 + bx + c = 0
\end{equation}

Його розв’язки залежать від дискримінанта $D$.

\begin{align}
D &= b^2 - 4ac \\
x_{1,2} &= \frac{-b \pm \sqrt{D}}{2a}
\end{align}
\section*{Проста таблиця}

\begin{tabular}{|c|c|c|}
\hline
1 & 2 & 3 \\
\hline
4 & 5 & 6 \\
\hline
7 & 8 & 9 \\
\hline
\end{tabular}
\section*{Таблиця з формулами}

\begin{tabular}{|c|c|c|}
\hline
$a^2$ & $b^2$ & $c^2$ \\
\hline
$\sqrt{a}$ & $\sqrt{b}$ & $\sqrt{c}$ \\
\hline
$\frac{1}{a}$ & $\frac{1}{b}$ & $\frac{1}{c}$ \\
\hline
\end{tabular}
\section*{Приклад розв'язання}

Розв’яжемо рівняння $2x^2 - 4x - 6 = 0$:

\begin{align*}
a &= 2,\quad b = -4,\quad c = -6 \\
D &= (-4)^2 - 4 \cdot 2 \cdot (-6) = 16 + 48 = 64 \\
x &= \frac{-(-4) \pm \sqrt{64}}{2 \cdot 2} \\
x &= \frac{4 \pm 8}{4} \\
x_1 &= 3,\quad x_2 = -1
\end{align*}

\end{document} % Кінець документа
